\documentclass[11pt]{article}
\usepackage{amsmath,amssymb,amsthm}
\usepackage[margin=1in]{geometry}
\usepackage[dvipsnames]{xcolor}

% hyperref should be the last package we load
\usepackage[pdftex,
colorlinks=true,
plainpages=false,
linkcolor=blue,
citecolor=Red,
urlcolor=black
]{hyperref}

\newtheorem{proposition}{Proposition}
\newtheorem{lemma}{Lemma}
\newtheorem{remark}{Remark}

\newcommand{\Om}{\Omega}
\newcommand{\dOm}{\partial\Omega}
\newcommand{\psilo}{\psi_\bot}
\newcommand{\psiup}{\psi^\top}
\newcommand{\psiloh}{\psi_{\bot,h}}
\newcommand{\psiuph}{\psi^\top_h}
\newcommand{\Cbot}{C_\bot}
\newcommand{\Ctop}{C^\top}
\newcommand{\uh}{u_h}
\newcommand{\Th}{\mathcal{T}_h}
\newcommand{\Vh}{\mathbb{V}_h}
\newcommand{\Vhr}{\mathring{\mathbb{V}}_h}
\newcommand{\Nh}{\mathcal{N}_h}
\newcommand{\Nhr}{\mathring{\mathcal{N}}_h}
\newcommand{\Sh}{\mathcal{S}_h}
\newcommand{\Ih}{I_h}
\newcommand{\Uh}{\mathcal{U}_h}
\newcommand{\Kh}{\mathcal{K}_h}
\newcommand{\Thbot}{\mathcal{T}_{h,\bot}}
\newcommand{\Thtop}{\Th^\top}
\newcommand{\Thplus}{\Th^+}
\newcommand{\Nhbot}{\mathcal{N}_{h,\bot}}
\newcommand{\Nhtop}{\Nh^\top}
\newcommand{\Lamh}{\Lambda_h}
\newcommand{\Lamhbot}{\Lambda_{h,\bot}}
\newcommand{\Lamhtop}{\Lambda_h^\top}

\title{Appendix: Pointwise a posteriori error control for box-constrained\\
elliptic variational inequalities}
\author{Ed Bueler \and Claude Sonnet 5 (Anthropic)}
\date{}

\begin{document}
\maketitle

\section{Setup}

Let $\Om$ be a bounded, polyhedral, not necessarily convex domain in $\mathbb{R}^d$
with $d \in \{1,2,3\}$. Let $f \in L^\infty(\Om)$ be a load function, let
$\psilo, \psiup \in H^1(\Om) \cap C^{0,\alpha}(\bar\Om)$, $0 < \alpha \le 1$, be
lower and upper obstacles satisfying the \emph{non-degeneracy assumption}
\begin{equation}
\label{eq:nondeg}
  \psilo < \psiup \qquad \text{in } \Om,
\end{equation}
and let
$g \in H^1(\Om) \cap C^{0,\alpha}(\bar\Om)$ be a Dirichlet boundary datum
satisfying the compatibility condition
\begin{equation}
\label{eq:compat}
  \psilo \le g \le \psiup \qquad \text{on } \dOm.
\end{equation}
Let $\mathcal{K}$ be the following non-empty, closed, convex, \emph{box-constrained}
subset of $H^1(\Om)$:
\begin{equation}
\label{eq:K}
  \mathcal{K} := \{v \in H^1(\Om) \mid \psilo \le v \le \psiup \text{ a.e.\ in } \Om,
    \ v = g \text{ on } \dOm\}.
\end{equation}
The continuous problem reads:
\begin{equation}
\label{eq:VI}
  u \in \mathcal{K}: \qquad
  \langle \nabla u, \nabla(v - u)\rangle \ge \langle f, v - u\rangle
  \qquad \text{for all } v \in \mathcal{K}.
\end{equation}
This is the immediate generalization of problem (1.1) of Nochetto, Siebert \&
Veeser (NSV03) \cite{NSV03} from a one-sided constraint $v \ge \chi$ to the
two-sided (\emph{bilateral}, \emph{box-constrained}) constraint
$\psilo \le v \le \psiup$; we state it in $(v-u)$ form throughout this
appendix (equivalent to NSV03's own $(u-v)$, $\le$ form by multiplying
through by $-1$), matching the convention used elsewhere in this literature
\cite{BF24}. As in NSV03, $\mathcal{K}$ admits a unique solution
$u$, which is also H\"older continuous.

\begin{remark}[Non-degeneracy]
The strict separation $\psilo < \psiup$ rules out \emph{pinched} points,
where the two obstacles touch and both constraints are simultaneously active.
This keeps the continuum (and per-triangle) analysis below free of a third,
degenerate case alongside lower- and upper-contact. A discrete
implementation must still cope with near-pinching, since active-set
membership there is decided against a numerical tolerance rather than exact
equality; that bookkeeping is a matter for the code, not this analysis.
\end{remark}

Let $\sigma \in H^{-1}(\Om) = \mathring H^1(\Om)^*$ be the functional
defined by
\begin{equation}
\label{eq:sigma}
  \langle \sigma, \varphi\rangle = \langle \nabla u, \nabla \varphi\rangle - \langle f, \varphi\rangle
  \qquad \text{for all } \varphi \in \mathring H^1(\Om).
\end{equation}

\begin{remark}[Sign convention]
\label{rem:signconvention}
This switches the sign relative to NSV03's formula (1.2); they define $\langle\sigma,\varphi\rangle = \langle f,\varphi\rangle
- \langle\nabla u,\nabla\varphi\rangle$, which forces $\sigma \le 0$
throughout in their unilateral (lower-obstacle) setting.  We use the
opposite sign so that $\sigma \ge 0$ on lower contact, and our sign choice matches the
traditional way a \emph{unilateral} complementarity condition is written:
$a \ge 0$, $b \ge 0$, $ab = 0$.   This is also
the convention already used by \texttt{viamr}'s \texttt{nsvmark()} and in \cite{BF24} \S 7.

The sign choice in \eqref{eq:sigma} does not by itself say anything about upper contact.  For a
unilateral \emph{upper}-obstacle problem ($v \le \psiup$, no lower
bound), for example, the same fixed formula \eqref{eq:sigma} gives $\sigma \le 0$ on
contact.  (Proof: test $v = u - t\varphi$, $\varphi \ge 0$, $t \ge 0$, the
free direction when only an upper bound is present.)

In any case, the reader should remember our sign convention as given by \eqref{eq:sigmasign} below.
\end{remark}

Unlike the unilateral case, $\sigma$ no longer has a fixed sign globally.
Write $\Cbot := \{u = \psilo\}$ and $\Ctop := \{u = \psiup\}$ for the closed
lower- and upper-contact sets; these are disjoint by the non-degeneracy
assumption. Then
\begin{equation}
\label{eq:sigmasign}
  \sigma \ge 0 \text{ on } \Cbot \qquad \text{and} \qquad \sigma \le 0 \text{ on } \Ctop,
\end{equation}
while $\sigma = 0$ on the open non-contact set $\{\psilo < u < \psiup\}$.

The discrete counterpart $\sigma_h$ appears below in \S 2.


\subsection{Separation of the contact sets}

The non-degeneracy assumption $\psilo < \psiup$ only says the two contact
sets $\Cbot, \Ctop$ are disjoint; it does not by itself say how far apart
they are. A mesh star touching \emph{both} contact sets would break the simple sign argument in \S 2 for the discrete residual $\sigma_h$ introduced below.

We need to control how much of the true contact set the discretization can
recognize as such.  Compactness of $\bar\Om$ turns the qualitative
separation into an explicit, data-determined one, quantifying that, as follows.

Let $[\psiup]_\alpha$ denote the H\"older seminorm of $\psiup$ on $\bar\Om$,
i.e.\ the least $L^\top \ge 0$ with $|\psiup(x) - \psiup(y)| \le L^\top
|x-y|^\alpha$ for all $x,y \in \bar\Om$.  This is finite since $\psiup \in
C^{0,\alpha}(\bar\Om)$, but note it can be zero, which motivates a convention below.

We \emph{assume}, as NSV03 does for the unilateral
problem, that the solution $u$ is also H\"older continuous on $\bar\Om$ with some exponent
$\beta \in (0,1]$ and seminorm $L_u \ge 0$.
%% TODO(citation): NSV03 attributes Hölder continuity of u, in the
%% unilateral case, to Frehse & Mosco (1982) [their ref. 7], not to
%% Rodrigues (1987) [their ref. 13, cited instead for existence/
%% uniqueness]. Neither has been read directly here; find and verify a
%% citation that actually covers the box-constrained/double-obstacle case
%% before this goes further, rather than asserting one.

\begin{proposition}[Separation of the contact sets]
\label{prop:separation}
Let
\begin{equation}
\label{eq:deltamin}
  \delta_{\min} := \min_{\bar\Om} (\psiup - \psilo) > 0,
\end{equation}
which is attained and positive by compactness of $\bar\Om$ and continuity
and non-degeneracy of $\psiup - \psilo$. If $\Cbot, \Ctop$ are both non-empty
then
\begin{equation}
\label{eq:r0}
  \operatorname{dist}(\Cbot, \Ctop) \ \ge\ r_0
  := \min\left\{ \left(\frac{\delta_{\min}}{2 L^\top}\right)^{1/\alpha},
                 \left(\frac{\delta_{\min}}{2 L_u}\right)^{1/\beta} \right\}.
\end{equation}
\end{proposition}

Regarding possible division by zero in \eqref{eq:r0}, we make the convention that a term is omitted from the minimum if its
denominator vanishes, i.e.\ if $\psiup$ or $u$ is constant.  That is, we regard the quotients as positive infinity in such cases.

\begin{proof}
Let $x \in \Cbot$, $y \in \Ctop$, and $r := |x-y|$. Since $u(x) = \psilo(x)$,
\begin{equation}
  \delta_{\min} \le \psiup(x) - \psilo(x) = \psiup(x) - u(x).
\end{equation}
Insert $\psiup(y) = u(y)$ and use the two H\"older bounds:
\begin{equation}
  \psiup(x) - u(x)
  = \bigl(\psiup(x) - \psiup(y)\bigr) + \bigl(u(y) - u(x)\bigr)
  \le L^\top r^\alpha + L_u r^\beta.
\end{equation}
Hence $\delta_{\min} \le L^\top r^\alpha + L_u r^\beta$. If $r < r_0$, then
$L^\top r^\alpha < \delta_{\min}/2$ and $L_u r^\beta < \delta_{\min}/2$ (each
by definition of $r_0$, dropping whichever term does not apply), so
$L^\top r^\alpha + L_u r^\beta < \delta_{\min}$, a contradiction. Thus
$r \ge r_0$, and since $x \in \Cbot$, $y \in \Ctop$ were arbitrary,
$\operatorname{dist}(\Cbot,\Ctop) \ge r_0$.
\end{proof}

Proposition~\ref{prop:separation} gives an explicit, data-determined
lower bound $r_0$ on the spatial separation of the two contact sets.  The bound depends on $\delta_{\min}$, on the H\"older constants of the obstacle data, and on the H\"older constant of the solution, which is itself controlled by $f,\psilo,\psiup,g$
through the cited regularity theory.  As it turns out (Remark~\ref{rem:starseparation2}), \S 2 does
\emph{not} need to impose $h < c\,r_0$ as a hypothesis for the estimator to
be well-defined and correct; $r_0$'s role is instead to bound the width of
the transition band $\Thplus$ once $h \lesssim r_0$, so that the
discretization's lower- and upper-contact classification tracks the true
contact sets closely, rather than under-recognizing them.


\section{Discretization}

Our notation follows NSV03 in most aspects.

Let $\{\Th\}$ be a sequence of conforming partitions of $\bar\Om$ into
closed simplices, namely closed intervals, closed triangles, or closed tetrahedra in dimensions $d=1,2,3$, respectively. Given a (closed) simplex $T$, $h_T$ is its diameter and $\rho_T$
is the maximal radius of a ball contained in $T$.  We assume shape-regularity, namely the bound
$\gamma(\Th) := \max_{T \in \Th} h_T/\rho_T \le \gamma^*$, with $\gamma^*$
independent of $h$.  Let $\Nh$ be the node set of $\Th$, $\Nhr \subset \Nh$
the interior nodes, and $\Sh$ the interior sides.

Let $\Vh$ be the space of
continuous piecewise-linear finite elements over $\Th$, with
$\Vhr := \Vh \cap H_0^1(\Om)$.  We denote its nodal basis by $(\varphi_h^z)_{z \in \Nh}$.

For $\omega
\subset \bar\Om$, let
\begin{equation}
\Uh(\omega) := \bigcup\{T \in \Th : T \cap \omega \ne \emptyset\}.
\end{equation}
This closed set is a (finite) union of the closed simplices
$T$.  Note that $T \cap \omega \ne \emptyset$ only requires $T$ to meet
$\omega$ in a single point, so $\Uh(T)$ for a simplex $T$ includes every simplex with a shared edge or face, but it even includes every simplex touching $T$ at a single shared vertex.  The star of a node $z \in \Nh$ is $\Uh(z) :=
\Uh(\{z\}) = \operatorname{supp} \varphi_h^z$, the closure of
$\{\varphi_h^z > 0\}$.

Let $\Ih$ be the Lagrange interpolation operator onto $\Vh$.  We write $\langle \cdot,\cdot\rangle_h$ for
the lumped $L_2$ scalar product,
\begin{equation}
\langle \varphi_h, \psi_h\rangle_h := \int_\Om \Ih(\varphi_h \psi_h)
\end{equation}
This is exactly as in NSV03.  Also as in that reference, in what follows $C$ denotes different constants
depending only on $\gamma^*$, $\Om$, and $d$, and $\preceq$ indicates $\le$
up to such a constant.


\subsection{Monotonicity of nodal interpolation, and $\Kh$}

\begin{lemma}[Order-preservation of $I_h$]
\label{lem:monotone}
If $\phi, \zeta \in C(\bar\Om)$ satisfy $\phi \le \zeta$ on $\bar\Om$, then
$\Ih \phi \le \Ih \zeta$ on $\bar\Om$, for any mesh $\Th$.
\end{lemma}
\begin{proof}
By linearity of $\Ih$, $\Ih \zeta - \Ih \phi = \Ih(\zeta - \phi)$, and
$\zeta - \phi \ge 0$ at every node. On each $T \in \Th$, $\Ih(\zeta-\phi)$
is affine and equals a convex combination (barycentric coordinates) of the
values of $\zeta - \phi$ at $T$'s vertices; a convex combination of
nonnegative numbers is nonnegative. Hence $\Ih(\zeta-\phi) \ge 0$ on $T$,
for every $T$.
\end{proof}

\begin{remark}
This is special to piecewise-\emph{linear} interpolation (a higher-order
Lagrange interpolant can overshoot and fail to preserve order); it needs no
acuteness or mesh-fineness assumption on $\Th$.
\end{remark}

Set $\psiloh := \Ih \psilo$ and $\psiuph := \Ih \psiup$. Applying
Lemma~\ref{lem:monotone} to $\psilo \le \psiup$ (non-degeneracy,
\eqref{eq:nondeg}) gives
\begin{equation}
\label{eq:disc-nondeg}
  \psiuph - \psiloh \ \ge\ \delta_{\min} \ >\ 0 \qquad \text{on } \bar\Om,
\end{equation}
with $\delta_{\min}$ exactly as in Proposition~\ref{prop:separation}; note
this holds for \emph{every} mesh, not just fine ones. Applying it again to
the compatibility condition $\psilo \le g \le \psiup$ on $\dOm$
\eqref{eq:compat} gives $\psiloh \le \Ih g \le \psiuph$ on $\dOm$. Hence, with
\begin{equation}
\label{eq:Kh}
  \Kh := \{v_h \in \Vh \mid \psiloh \le v_h \le \psiuph \text{ in } \Om,
    \ v_h = \Ih g \text{ on } \dOm\},
\end{equation}
$\Kh$ is non-empty, closed, and convex for \emph{every} mesh $\Th$.  As in NSV03, $\Kh$ is in general not a subset of
$\mathcal{K}$ defined by \eqref{eq:K}, so in this sense our approximation is non-conforming.

The discrete problem corresponding to \eqref{eq:VI} reads
\begin{equation}
\label{eq:discVI}
  u_h \in \Kh: \qquad
  \langle \nabla u_h, \nabla(v_h - u_h)\rangle \ge \langle f, v_h - u_h\rangle
  \qquad \text{for all } v_h \in \Kh,
\end{equation}
which admits a unique solution.

\subsection{The discrete residual and its sign}

For interior nodes $z \in \Nhr$, define $\sigma_h \in \Vh$ nodally by
\begin{equation}
\label{eq:sigmah}
  \langle \sigma_h, \varphi_h^z \rangle_h
  = \langle \nabla u_h, \nabla \varphi_h^z \rangle - \langle f, \varphi_h^z \rangle,
\end{equation}
whence $\sigma_h(z) = \bigl(\int_\Om \varphi_h^z\bigr)^{-1}
\bigl(\langle \nabla u_h, \nabla \varphi_h^z\rangle - \langle f, \varphi_h^z\rangle\bigr)$.

\begin{remark}[Consistent with $\sigma$]
This is exactly NSV03's own formula (2.1), with the same sign flip as
$\sigma$ itself (Remark~\ref{rem:signconvention}) --- $\sigma_h$ is not a
second, independent sign choice, it is the discrete instance of the same
$\sigma$ already fixed in \S 1. Concretely, \eqref{eq:sigmah} is the
Galerkin (weak-form) instance of the strong-form nodal residual underlying
BF24's mixed complementarity classification \cite{BF24} \S 7, and matches
the convention already used by \texttt{viamr}'s \texttt{nsvmark()}.
\end{remark}

Let $\Nhbot := \{z \in \Nh : u_h(z) = \psiloh(z)\}$ and $\Nhtop := \{z \in
\Nh : u_h(z) = \psiuph(z)\}$ be the discrete lower- and upper-contact node
sets; these are disjoint by \eqref{eq:disc-nondeg}. Split $\Th$, exactly as
NSV03 splits into $\Th^0, \Th^+$, but now into three parts:
\begin{align}
\label{eq:Thsplit}
  \Thbot &:= \{T \in \Th : \Uh(T) \subset \{u_h = \psiloh\}\}, \\
  \Thtop &:= \{T \in \Th : \Uh(T) \subset \{u_h = \psiuph\}\}, \\
  \Thplus &:= \Th \setminus (\Thbot \cup \Thtop).
\end{align}
$\Lamhbot$ and $\Lamhtop$ denote the corresponding unions of elements
(discrete lower- and upper-contact sets).

\begin{remark}[These are \texttt{viamr}'s thin active sets]
\label{rem:thinactive}
$\Thbot$ and $\Thtop$ are exactly what \texttt{viamr} calls
\texttt{thinelemactive()}: an element is counted only if it \emph{and every
element sharing a node with it} is active, i.e.\ the naive element-wise
active set eroded by one ring. That erosion is precisely the content of
$\Uh(T) \subset \{u_h = \psiloh\}$ --- $\Uh(T)$ \emph{is} $T$'s ring of
neighbors (every $T' \in \Th$ meeting $T$ in a node), so requiring the whole
ring to be active is the same condition, just written in NSV03's star
notation instead of \texttt{viamr}'s dilation-based one. Concretely,
$\Thbot = \texttt{thinelemactive(u\_h, psilo\_h, boxside="lower")}$ and
$\Thtop = \texttt{thinelemactive(u\_h, psiup\_h, boxside="upper")}$, using
the \texttt{boxside} generalization already in \texttt{viamr}. This is not
a coincidence: NSV03 needed the one-ring erosion for exactly the same
reason \texttt{thinelemactive()} was built for AMR marking --- a single
active element adjacent to the free boundary is not yet safely in the
interior of the contact set, so testing (or refining, or excluding from
refinement) based on its value alone is unsound; only once a whole
neighborhood agrees is the sign (or safety) of that neighborhood-based
decision unconditional. Proposition~\ref{prop:nodalsign} is the a
posteriori-estimator-side reason \texttt{thinelemactive()}'s erosion is the
right notion of ``safely interior,'' not merely a heuristic.
\end{remark}

\begin{lemma}[Vanishing of an affine function on a simplex]
\label{lem:affine}
Let $w$ be affine on a simplex $T$ with $w \ge 0$ at every vertex of $T$.
Then $w = 0$ somewhere in $T$ if and only if $w = 0$ at some vertex of $T$.
\end{lemma}
\begin{proof}
($\Leftarrow$) is trivial. For ($\Rightarrow$), write a point $x \in T$ in
barycentric coordinates, $x = \sum_i \lambda_i x_i$, $\lambda_i \ge 0$,
$\sum_i \lambda_i = 1$, so $w(x) = \sum_i \lambda_i w(x_i)$. If $w(x_i) > 0$
for every vertex $x_i$, then $w(x) > 0$ for every $x \in T$ (some $\lambda_i
> 0$ always, and all terms are non-negative with at least one strictly
positive). Contrapositive gives the claim.
\end{proof}

\begin{proposition}[Nodal sign, wholly-contact case]
\label{prop:nodalsign}
If $z \in \Nhr$ and $\Uh(z) \subset \{u_h = \psiloh\}$ (i.e.\ $z$'s star is
wholly lower-contact), then $\sigma_h(z) \ge 0$. Symmetrically, if $\Uh(z)
\subset \{u_h = \psiuph\}$, then $\sigma_h(z) \le 0$.
\end{proposition}
\begin{proof}
Suppose $\Uh(z) \subset \{u_h = \psiloh\}$. Since $\psiuph - \psiloh \ge
\delta_{\min} > 0$ on $\bar\Om$ \eqref{eq:disc-nondeg} and $u_h = \psiloh$
throughout $\Uh(z)$, also $\psiuph - u_h \ge \delta_{\min} > 0$ throughout
$\Uh(z)$; in particular $\Uh(z) \cap \Nhtop = \emptyset$. Fix $t \in [0,
\delta_{\min}/\max_{\Uh(z)} \varphi_h^z]$ and set $v_h := u_h + t
\varphi_h^z \in \Vh$. Since $\varphi_h^z \ge 0$ and $u_h = \psiloh$ on
$\Uh(z)$, $v_h \ge \psiloh$ there (and $v_h = u_h \ge \psiloh$ elsewhere);
by the choice of $t$, $v_h \le u_h + \delta_{\min} \le \psiuph$ on $\Uh(z)$
too (and $v_h = u_h \le \psiuph$ elsewhere). So $v_h \in \Kh$. Testing
\eqref{eq:discVI} with this $v_h$ and dividing by $t \ge 0$ gives $\langle
\nabla u_h, \nabla \varphi_h^z\rangle \ge \langle f, \varphi_h^z\rangle$,
i.e.\ $\sigma_h(z) \ge 0$ by \eqref{eq:sigmah}. The upper case is symmetric,
testing with $t \le 0$.
\end{proof}

\begin{remark}[Why no separation hypothesis is needed here]
\label{rem:starseparation2}
This is exactly the case NSV03's own estimator uses $\sigma_h$'s sign for:
elements of $\Thbot$ or $\Thtop$. For a node whose star meets \emph{both}
contact types (necessarily an element of $\Thplus$, by
\eqref{eq:Thsplit}), the perturbation argument above breaks --- moving
$t>0$ toward the free lower direction can run directly into an upper-contact
point elsewhere in the star, for any $t>0$, however small (Lemma
\ref{lem:affine} with the roles reversed) --- but $\Thplus$ elements never
call on $\sigma_h$'s sign in the estimator to begin with (see the
forthcoming efficiency/lower-bounds section), so no separation hypothesis
is actually load-bearing for correctness.
Proposition~\ref{prop:separation}'s $r_0$ still matters for
\emph{sharpness}: once $h \lesssim r_0$, $\Thplus$ shrinks to a thin
transition band tracking the true free boundary, rather than swallowing up
contact elements that should have been classified $\Thbot$ or $\Thtop$.
\end{remark}

%% TODO: remaining sections, mirroring NSV03's structure, not yet written:
%%
%%  (a) Boundary formula for sigma_h at Dirichlet nodes z in N_h \ N_h^o: no
%%      perturbation freedom there (v_h = I_h g is an equality constraint),
%%      so sigma_h(z) there is a diagnostic quantity, not derived from
%%      stationarity as in Proposition~\ref{prop:nodalsign}. Need the
%%      two-sided analogue of NSV03's boundary-node formula (their page 169):
%%      a symmetric (R -+ bnd)^{+/-} truncation construction, worked out in
%%      discussion 2026-08-18 but not yet transcribed into this file.
%%  (b) Positivity-preserving interpolation Pi_h (NSV03 sec. 2.2, the
%%      Chen-Nochetto ball-averaging operator): restate; properties (2.5)-
%%      (2.8) should carry over unchanged, since Pi_h does not depend on
%%      obstacle shape.
%%  (c) Galerkin functional G_h (NSV03 sec. 3): generalize the T_h^0/T_h^+
%%      split there to Thbot/Thtop/Thplus as in section 2 above; check
%%      whether NSV03's almost-Galerkin-orthogonality (3.4) and
%%      representation formula (3.5) need a second correction term for the
%%      second obstacle.
%%  (d) Barriers (NSV03 sec. 4): upper/lower barrier construction via the
%%      Riesz representation w of G_h; the continuous-maximum-principle
%%      argument looks obstacle-shape-agnostic and should carry over
%%      directly, but has not been checked.
%%  (e) Pointwise estimates (NSV03 sec. 5): Green's function bound and
%%      Holder continuity of w; likely unchanged from NSV03 modulo notation.
%%  (f) Efficiency / lower bounds (NSV03 sec. 6): Prop. 6.2/6.3 analogues
%%      (bubble functions), needed for both Thbot and Thtop.
%%  (g) Still open: citation for Holder continuity of u in the
%%      box-constrained case (see TODO(citation) in section 1 above).

\bibliographystyle{plain}
\begin{thebibliography}{9}
\bibitem{NSV03}
R.~H.~Nochetto, K.~G.~Siebert, A.~Veeser.
\emph{Pointwise a posteriori error control for elliptic obstacle problems}.
Numer. Math. \textbf{95} (2003), 163--195.
\bibitem{BF24}
E.~Bueler, P.~E.~Farrell.
\emph{A full approximation scheme multilevel method for nonlinear
variational inequalities}.
SIAM J. Sci. Comput. \textbf{46} (2024), A2421--A2444.
\end{thebibliography}

\end{document}
